\section{Source and evidence map}
\label{app:artifacts}

The source snapshot for this manuscript's established protocol and experiment
records is commit
\texttt{228c944fc476e3e57d6c08063cacb19bca5b360e}
of \href{https://github.com/billlza/q-periapt}{\texttt{billlza/q-periapt}}.
Each retained execution keeps its original source and tool identity.
The local record model and accounting proofs in this manuscript are mathematical
arguments; the following entries identify the independently recorded
mechanization and native observations.

\begin{longtable}{@{}>{\raggedright\arraybackslash}p{0.25\linewidth}>{\raggedright\arraybackslash}p{0.69\linewidth}@{}}
\toprule
Result & Reproduction source\\
\midrule
\endhead
Joint symbolic replay &
\href{https://github.com/billlza/q-periapt/tree/228c944fc476e3e57d6c08063cacb19bca5b360e/research/sdk-alpha1/evidence/20261104-multi-epoch-replay}{20261104-multi-epoch-replay}.
The original archive is \texttt{CAPTURE.zip}; the checked certificate member is
\texttt{replay-01/01-baseline-replay.proof.spthy}.\\[5pt]
Positive two-round trace &
\href{https://github.com/billlza/q-periapt/tree/228c944fc476e3e57d6c08063cacb19bca5b360e/research/sdk-alpha1/evidence/20261105-two-round-ground-trace}{20261105-two-round-ground-trace}.\\[5pt]
Delayed-receipt ground trace &
\href{https://github.com/billlza/q-periapt/tree/228c944fc476e3e57d6c08063cacb19bca5b360e/research/sdk-alpha1/evidence/20261105-delayed-receipt-ground-trace}{20261105-delayed-receipt-ground-trace}.\\[5pt]
Native reservation disclosure &
\href{https://github.com/billlza/q-periapt/blob/228c944fc476e3e57d6c08063cacb19bca5b360e/research/continuity-identity-candidate/RESERVATION_DISCLOSURE.md}{RESERVATION\_DISCLOSURE.md}
and \path{src/durable/messages/tests/reservation_disclosure.rs}
in the identity candidate.\\[5pt]
Aggregate-copy execution &
\href{https://github.com/billlza/q-periapt/tree/228c944fc476e3e57d6c08063cacb19bca5b360e/research/sdk-alpha1/evidence/20261105-aggregate-reservation-copy}{20261105-aggregate-reservation-copy}.\\[5pt]
Root, traffic and phase mapping &
\href{https://github.com/billlza/q-periapt/blob/228c944fc476e3e57d6c08063cacb19bca5b360e/docs/continuity/RECOVERY_DEPENDENCIES_V1.md}{RECOVERY\_DEPENDENCIES\_V1.md};
native \texttt{rekey/response.rs} and \texttt{rekey/completion.rs}.\\[5pt]
SPQR snapshot corpus &
\href{https://github.com/billlza/q-periapt/tree/228c944fc476e3e57d6c08063cacb19bca5b360e/research/continuity-spqr-reference}{continuity-spqr-reference}:
\texttt{TRACE\_CORPUS.json}, \texttt{COMPROMISE\_CORPUS.json}, and
\texttt{COMPROMISE\_EXPERIMENT.md}.\\[5pt]
Whole-KEM corpus &
\href{https://github.com/billlza/q-periapt/tree/228c944fc476e3e57d6c08063cacb19bca5b360e/research/continuity-whole-kem-reference}{continuity-whole-kem-reference}:
\texttt{TRACE\_CORPUS.json}, \texttt{README.md}, and
\texttt{ACTIVE\_COMPROMISE.md}.\\[5pt]
Computational component arguments &
\href{https://github.com/billlza/q-periapt/blob/228c944fc476e3e57d6c08063cacb19bca5b360e/docs/continuity/COMPUTATIONAL_RECOVERY_V1.md}{COMPUTATIONAL\_RECOVERY\_V1.md}
and its linked KEM adapter, hidden-record, shared-HMAC,
root-extraction and adaptive-target notes.\\
\bottomrule
\end{longtable}

The complete seven-schedule reference comparison is reproduced by
\path{artifacts/reference-comparison/analyze_reference.py} with
\texttt{--check}. Its four retained corpus files are bound to the source
snapshot above by Git blob and SHA-256 identities in
\path{artifacts/reference-comparison/source_identity.json}.
\path{analysis.json} in that directory records all 420 snapshot cases and
the 35 configuration--schedule rows; the generated
\path{reference-schedule-table.tex} supplies Table~\ref{tab:reference-schedules}.
The independent review checks every displayed count and decimal mean using
separate integer arithmetic and rejects eight altered-input controls.

The jointly checked certificate has SHA-256
\begin{center}
\small\texttt{dde5d06a071b1560da0d4c28c3f1fd81a3c39708f9fd8564ee391ddc362bcf7f}.
\end{center}
Its source-base receipt identifies
\texttt{6084f3bfb09333f6b66291b38809ed512296a27b}.
The fixed 4 November replay retained all 40 queries and its exact tool,
resource and output receipts. Reading and checking this retained artifact while
preparing the manuscript is distinct from executing another Tamarin run.



\subsection{Mechanization and replay record}

The theorem is the retained query
\path{multi_epoch_fresh_reservations_recover_symbolic_message_secrecy}.
Its certificate was replayed with all 24 reusable dependencies in the same
Tamarin process. The independent 4 November replay used Tamarin 1.12.0 and
Maude 3.5.1, one worker, a 12-GiB heap bound and the unchanged 900-second limit;
it completed in 816.903 seconds. The output certificate was byte-identical to
the retained input. Appendix~\ref{app:artifacts} gives its digest and archive
member, making the theorem's exact formula and proof available for inspection.

The complete inventory contains 40 queries: 27 verified and 13 incomplete.
The latter comprise eleven multi-round existence queries and two historical
all-trace queries outside the main theorem's dependency group. Two subsequent
finite ground-trace constructions supply the concrete executions discussed
below. Their independently checked formulas and original certificate inventory
remain separately identifiable.


\subsection{Guarded-normalization supplement}

The complete manuscript extends the previously retained source at commit
\begin{center}
\small\texttt{1c68aa666db22d36ee136be1ec712545a9e74f03}
\end{center}
under \path{paper/continuity}. The manuscript and its directly used evidence
are maintained in the repository's
\href{https://github.com/billlza/q-periapt/pull/119}{continuous manuscript PR}.
The filenames below are relative to \path{paper/continuity} in that distribution.

The explicit normalization theorem and its proof are in
\path{guarded-provenance.tex}; the conjunctive-exposure proposition and proof
are in \path{acquisition-nonadditivity.tex}. The executable supplement is in
\path{artifacts/guarded-provenance/}:
\begin{itemize}
\item \path{check_exposure.py} evaluates both record semantics;
\path{results.json} records the finite enumeration and history checks.
\item \path{review_exposure.py} checks 17 independently stated expected
results; \path{review_results.json} records their outcomes.
\item \path{review.md} documents the theorem and implementation review;
\path{review_optimized_guard.json} records fail-closed optimized-mode
rejection. The separately named
\path{review_optimized_run_before_guard.json} retains the earlier execution
that motivated the guard.
\item \path{native_source_map.json} records source correspondence for the
journal-to-wrapping-key derivation, durable image serialization, pending
targets and the retained aggregate-copy test. It pins five source files
at commit \texttt{228c944}; this inspection records no new native execution.
\end{itemize}
From that artifact directory, the reproduction commands are
\begin{quote}
\small
\begin{verbatim}
python3 check_exposure.py --output results.json
python3 review_exposure.py
\end{verbatim}
\end{quote}
These commands exercise the finite record model and its independent expected
cases. The retained native and Tamarin evidence remains at the source
checkpoints listed in the preceding map.

\subsection{Constructive image/intent simulation}
\label{app:opaque-artifacts}

Section~\ref{sec:opaque-rendering} integrates the complete finite-interface
algorithm and proof from \path{artifacts/record-rendering/opaque-renderer-v1.2.md}.
The corresponding mathematical review is in
\path{artifacts/record-rendering/independent-review.md}. The retained source
modules defining the commitment experiment are
\href{https://github.com/billlza/q-periapt/blob/228c944fc476e3e57d6c08063cacb19bca5b360e/docs/continuity/HIDDEN_RECORD_COMMITMENTS_V1.md}{\nolinkurl{HIDDEN_RECORD_COMMITMENTS_V1.md}}
and
\href{https://github.com/billlza/q-periapt/blob/228c944fc476e3e57d6c08063cacb19bca5b360e/docs/continuity/MULTI_RECORD_COMMITMENTS_V1.md}{\nolinkurl{MULTI_RECORD_COMMITMENTS_V1.md}}.
Their actual domains and complete oracle budgets define the
$\epsilon_{\rm com}$ term; the new table construction supplies the
image/intent implementation after that replacement.

The source correspondence of Section~\ref{sec:native-correspondence} is
recorded in \path{artifacts/native-interface/native_interface_map.md} and
\path{artifacts/native-interface/map.json}. The machine-readable map binds
the source ranges and file identities to the fixed protocol checkpoint.
\path{artifacts/native-interface/native_map_validation.json} records the
source checks. These are source observations associated with the declared
interfaces; historical execution receipts retain their own dates and
source identities.

\subsection{Immutable report composition}
\label{app:report-artifacts}

Section~\ref{sec:report-composition} gives the finite-interface algorithm,
reverse-order theorem and real-disclosure corollary. Its complete proof note
is \path{artifacts/report-composition/multi_report_theorem.md}; the independent
review is \path{artifacts/report-composition/multi_report_review.md}.
The fixed-source
\href{https://github.com/billlza/q-periapt/blob/228c944fc476e3e57d6c08063cacb19bca5b360e/docs/continuity/CLEANUP_REPORT_MAC_V1.md}{\nolinkurl{CLEANUP_REPORT_MAC_V1.md}}
records the fixed-source one-row interface. The paper states its operational
game in Section~\ref{sec:report-composition} and its primitive instantiation
in Appendix~\ref{app:primitive-inputs}.
\href{https://github.com/billlza/q-periapt/blob/228c944fc476e3e57d6c08063cacb19bca5b360e/docs/continuity/SHARED_ROOT_HMAC_V1.md}{\nolinkurl{SHARED_ROOT_HMAC_V1.md}}
specifies the compression-oracle functions and their query regimes.
In the paper's notation, $\Phi_{256}$ bounds the two secret HMAC key-pad
replacements, and $E_{\rm NMAC}$ bounds the ensuing row replacement with
the stated length and query budgets. The composition theorem accepts the
one-row distinguishing bound as a modular input and adds the complete
ordinary-query work of every derived reduction.

\subsection{Candidate-local computational connection}

Section~\ref{sec:candidate-local-accounting} incorporates the frozen proof
fragment paired with \path{artifacts/candidate-local/candidate_local_composition.md}.
The independent derivation and its artifact identities are
recorded in \path{artifacts/candidate-local/independent_review.md} and
\path{artifacts/candidate-local/independent_review.json}. They check the
all-continuations stopping certificate, the exact native-score identity,
the selected-candidate residual, and the local component comparisons.
The score identity applies to bounded scores; the positive/negative-event
conversion is specialized to the message experiment's three-valued score.

\subsection{Joint schedule and live recovery profile}
\label{app:joint-artifacts}

Section~\ref{sec:joint-continuation} and Appendix~\ref{app:joint-primitives}
give the joint schedule and its shared-oracle primitive derivations.
The complete proof record is
\path{artifacts/joint-continuation/joint_continuation.md}; its separate
review and frozen source identities are retained in the same directory.
The live recovery profile in \path{artifacts/live-profile/source_profile.md}
maps the actual root, confirmation, traffic, message and ACK calls, together
with the private application and persistent-record interfaces. The map and
encoding inventory record the exact full information strings and source
locations. Its 24 source files are byte-identical between the protocol pin
\sourcecommit{} and the inspected development revision
\texttt{1c68aa66}; this identity check is separate from the retained native
execution receipts.

The live-core instantiation and joint authentication consequence are retained
in \path{artifacts/confirmation-traffic/}. Their proof and separate reviews
record the complete observation surface, the reservation-selector witness,
the native confirmation cores and the original encryption ledger. The source
map fixes each inspected file at the protocol checkpoint. The authentication
argument starts at the uniform-root interface and assigns its own terminal
event to any preceding comparison.

\subsection{ContextBound and original-operation inputs}
\label{app:contextbound-source}

The core \texttt{combine\_policy\_bound} and \texttt{combine\_validated}
functions in
\href{https://github.com/billlza/q-periapt/blob/228c944fc476e3e57d6c08063cacb19bca5b360e/crates/q-periapt-core/src/lib.rs}{\path{crates/q-periapt-core/src/lib.rs}}
fix the length-delimited component and context encoding. The SDK replay
constructor's \texttt{reserve}, \texttt{open} and \texttt{binding} operations in
\href{https://github.com/billlza/q-periapt/blob/228c944fc476e3e57d6c08063cacb19bca5b360e/crates/q-periapt-sdk/src/expert/replay.rs}{\path{crates/q-periapt-sdk/src/expert/replay.rs}}
retain the admitted operation and its original tape. The rekey response
module's \texttt{kem\_context} in
\href{https://github.com/billlza/q-periapt/blob/228c944fc476e3e57d6c08063cacb19bca5b360e/research/continuity-identity-candidate/src/durable/messages/rekey/response.rs}{\path{research/continuity-identity-candidate/src/durable/messages/rekey/response.rs}}
binds the complete retained offer into the application context.
The artifact also retains an earlier Q-Periapt manuscript on ContextBound
binding and implementation assurance. Those materials document the reused
cryptographic substrate; the original reservation, acquisition and persistent
simulation results are developed in the present paper.

The connection artifact, \path{artifacts/contextbound-connection/},
contains the complete proof, independent review and source induction for the
next rekey of one established session. Its source map fixes fourteen files
at \sourcecommit{} and the inspected development checkpoint. The
self-contained checker recomputes six retained core vectors and exercises
complete-tuple aliases, original-operation retries and component-challenge
routing, with explicit incorrect-cache controls. The source audit is a
separate command requiring the declared native revisions; its receipt records
source identity and static assertions.

\paragraph{Relation to earlier project material.}
Earlier project work on SkyBridge Compass studies policy-governed
authenticated establishment, complete-transcript binding, and a control-only
recovery phase before application traffic
(\href{https://github.com/billlza/Skybridge-Compass/blob/d49d2a9ed4e1fb00905adeb06c02989562e375af/Docs/TDSC-2026-01-0318_IEEE_Paper_SkyBridge_Compass_patched.tex}{unpublished manuscript, revision \texttt{d49d2a9}}).
It supplies the operation-admission and traffic-release perspective. Here we
follow the admitted operation's original randomness through persistence and
acquisition, and develop record and report simulations. The artifact's
\path{notes/skybridge-scope.md} fixes the source versions and allocates material
between these research questions.


\subsection{Finite multiple-root entry}

Section~\ref{sec:multiple-root-entry} gives the finite controller, joint
root-entry comparison and first-message corollary. The complete argument,
including the class-allocation reduction, query budgets and a constructive
controller, is in
\path{artifacts/multiple-root-entry/multiple_root_entry_proof.md}.
The same directory contains the frozen TeX module, independent derivation,
25-point mathematical review and exact file identities. The controller
permits available unmarked roots to become later salts; the marked root and
all its aliases use the retained schedule handle. The review checks this
interface and its composition with the preceding primitive bounds.

The prose correspondence record in \path{artifacts/editorial-review/}
links the current presentation to the retained mathematical inputs. Historical
proof snapshots keep their exact identities; the semantic review checks the
current definitions, hypotheses, interfaces and budgets against them.

\subsection{SDK purpose-key correspondence}

Section~\ref{sec:sdk-purpose-connection} uses the owned purpose-key API in
\href{https://github.com/billlza/q-periapt/blob/228c944fc476e3e57d6c08063cacb19bca5b360e/crates/q-periapt-sdk/src/purpose.rs}{\path{crates/q-periapt-sdk/src/purpose.rs}}.
Its complete proof and source map are retained in
\path{artifacts/sdk-purpose-connection/}. The map checks the effective
information fields, HKDF calls, owner guards and retained native tests at
the protocol pin and development checkpoint. The accompanying review
checks the variable-length compression budgets, complete-input aliases,
owner transitions and selected-key experiment.

The Rust check job at development revision \texttt{1c68aa66} passed all
five purpose-key tests on 7 November 2026. The retained receipt lists their
outcomes, covering the RFC~5869 vector, domain separation, dual-peer
derivation, runtime closure, input admission and trust-root binding.
The source correspondence and mathematical review record the additional
manuscript claims.

\subsection{Executed sources and signing-reservation interfaces}
\label{app:source-attribution}

Section~\ref{sec:executed-source-attribution} has its complete proof,
source identities and independent mathematical review in
\path{artifacts/source-attribution/}. The directly inspected native
revision is \texttt{72acaaccde1ca9ea27e5af022da5fbce932ed4e9}; five of its
seven inspected files additionally match the earlier native-pin receipt.
The source map records the exact reservation/persistence/computation order,
complete-wire bindings, original-signing-randomizer encoding and the
separate signer-seed-file interface.

The executable model explores 34 recovered states for generation and
34 for encapsulation, with up to three original draws and two epochs.
Its 486 call-boundary edges include failed writes with absent or retained
targets, private-computation failure, process loss, release denial and
completion. Six deliberate source-tracking mutations have explicit
counterexamples. An independent expected matrix covers offer, response
and final signature aliases; the response alias permits proposer traffic
allocation with unchanged sources and changed complete traffic information.
The signature checks preserve original randomizers and private-query
ledgers across unpublished results, and enumerate the generic separation's
64 small key/randomizer assignments. These are finite interface checks
and exact source audits. Their receipts keep the native Rust and
mechanized-proof execution status separate.

\subsection{Reservation-to-ordinary signature connection}
\label{app:signature-bridge}

Section~\ref{sec:reservation-ordinary-bridge} has its full reduction,
source identities, finite controls and independent internal argument review
in \path{artifacts/signature-bridge/}. The native correspondence uses
revision \texttt{358b921ab3469ce49f166633fc9a762ff6ed7e3d} and the
lockfile-authenticated \texttt{fips204} 0.4.6 source archive. Its upstream
revision is \texttt{8d456d989dce681233ef963d3c3306da8f76d7c7}.
The retained packet includes the inspected sources and their licenses, so
its identity and call-length checks run from the manuscript checkout.

The controls distinguish exact pair caching from per-slot oracle queries,
preserve original-message freshness under descriptor aliases, and record
private invocations whose publication fails. Independent finite enumeration
checks the shared-XOF construction's joint distribution and its disjoint
hidden-input sectors. Owner controls distinguish declared lineage, actual
signing-state identity and public-key equality, and exercise non-target
disclosure. The mathematical review covers the complete two-comparison
ROM/QROM argument and its ordinary native signing endpoint; executable
controls validate the specified finite distributions and source interfaces.

\subsection{Native authentication composition}
\label{app:native-authentication}

Section~\ref{sec:native-authentication} gives the complete selected-root
authentication proof. Its argument review and executable accounting
controls are retained in \path{artifacts/native-authentication/}.
The controls check the native Boolean decomposition on 160 compatible
assignments, including free and charged stops. Separate finite examples
distinguish the selected authentication event from the confidentiality
score events, joint eligibility from conditional probability, and the
selected root from other root histories. The proof uses the source and
primitive correspondences identified in the preceding supplements; its
event-bit comparisons retain their own complete continuation budgets.

\subsection{Matched directed-opportunity executions}
\label{app:matched-comparison}

Section~\ref{sec:matched-opportunities} uses research revision
\texttt{fec8472060728cc1bb35cf15042c117dda00e336}, executed by GitHub run
\href{https://github.com/billlza/q-periapt/actions/runs/42537214047}{42537214047}.
Its checkout is the synthetic merge \texttt{0d12e77d}; preparation archives
protocol source \texttt{358b921a} into an isolated execution tree. The retained
preparation and executable receipts identify the native driver, SPQR and its
64-byte chunk variant, and the whole-KEM reference. The run completed all
49 profile--schedule pairs.

\path{artifacts/matched-comparison/} retains the complete executed report,
its 49-row JSON/CSV projection, original public-audit and qualification
receipts, source digests and deterministic manuscript-table generator.
The compressed report preserves the exact original JSON bytes. The generator
checks complete profile/schedule and capture-owner inventories, admission and
exposure partitions, frame budgets, full-envelope bytes and the distinction
between native physical persistence and reference state returns. Input
controls exercise altered denominators, duplicate rows, inconsistent request
sets, fabricated physical counters and omitted frame-envelope work.

The original public audit recomputed 14 native traces, 35 reference
trace/snapshot cases and 35 public envelope schedules, and checked all 126
public-evidence manifest entries. The hosted run performed the real native
plaintext and reference AEAD key checks. The retained qualification receipt
binds archive 13641008832, its SHA-256, the exact source and execution, and
these verification outcomes. The manuscript projection preserves those
execution identities and can be regenerated from the retained packet alone:
\begin{samepage}
\begin{quote}
\small
\begin{verbatim}
python3 artifacts/matched-comparison/analyze_matched.py --check
python3 artifacts/matched-comparison/check_input_controls.py
\end{verbatim}
\end{quote}
\end{samepage}

\subsection{Retry-readiness ablation}
\label{app:retry-readiness}

Section~\ref{sec:retry-readiness} uses research revision
\texttt{29f1c362b67bcd06123ea1bef847af82e9f9e100}, executed by GitHub run
\href{https://github.com/billlza/q-periapt/actions/runs/42551086845}{42551086845}
at synthetic checkout \texttt{23f600f5}. The fixed protocol source remains
\texttt{358b921a}. Its source comparison covers 2,687 files and identifies
only the addition of the two experiment modules to the native test file;
the original immediate-policy driver retains its exact bytes. The executed
test binary is hash-bound before and after both policy runs.

\path{artifacts/retry-readiness/} retains the exact compressed execution
report, source and executable identities, fourteen paired result rows, the
source/readiness control log and the independent public-trace audit. The
source-bound provenance receipt identifies job 128946601334 and archive
13650034930, checks all 42 archive members, and records fourteen passing
source/readiness controls. The trace audit checks 29 public evidence files,
recounts all 28 executions, and matches all fourteen fresh immediate-policy
rows to the retained matched comparison.

The table generator checks the two policies' complete profile/schedule
inventories, application and exposure partitions, emitted-frame classes and
bytes, public retry ages and physical persistence counters. It compares the
executed report with the independent recount and the fourteen-pair result
projection before emitting the manuscript table. The complete archived
counts distinguish first commitments from deliveries and retain each
post-capture denominator. Reproduction from the compact packet is:
\begin{quote}
\small
\begin{verbatim}
python3 artifacts/retry-readiness/analyze_retry.py --check
\end{verbatim}
\end{quote}
